% TOP_PROBLEM: {"id": 2943, "title": "Kirby Problem 4.67", "category_id": 7, "category": {"id": 7, "name": "topology", "display_name": "Topology", "description": "Properties preserved under continuous deformations.", "slug": "topology", "order_index": 7, "created_at": "2026-07-31T15:26:25.671Z"}, "status": "open"}
% FIRST_POSTED: null
% ENTRY_KIND: statement_only
% Imported from: batch13.tex; UnsolvedMath ID 2943
\documentclass[11pt]{article}
\usepackage[margin=25mm]{geometry}
\usepackage{fontspec}
\setmainfont{Latin Modern Roman}
\usepackage{amsmath,amssymb,mathtools}
\usepackage{unicode-math}
\setmathfont{Latin Modern Math}
\usepackage{xurl,hyperref}
\hypersetup{colorlinks=true,urlcolor=blue}
\setcounter{secnumdepth}{0}
\setlength{\parindent}{0pt}
\setlength{\parskip}{5pt}
\setlength{\emergencystretch}{3em}
\newcommand{\sref}[2]{\hyperlink{source:#1:#2}{[S#2]}}
\newcommand{\cataloguescope}[1]{\paragraph{Recorded target.}#1}
\begin{document}
% UnsolvedMath ID 2943; source catalogue code KP-4.67
\setcounter{section}{2942}
\section{Kirby Problem 4.67}\label{2943}
{\small\sffamily UnsolvedMath ID 2943 · Topology}
\subsection{Definitions and mathematical statement}

\paragraph{Shared notation.}$\mathbb N=\{1,2,\ldots\}$, and $\mathbb Z,\mathbb Q,\mathbb R,\mathbb C$ denote the integers, rationals, reals and complex numbers. Any other conventions are specified locally.


Let $\operatorname{Diff}^+(S^4)$ be the smooth-topology group of orientation-preserving diffeomorphisms of the standard four-sphere. Its component group $\pi_0$ is the group of smooth isotopy classes under composition.

The four-dimensional hollow barbell $B$ is a slight regular neighborhood in $\mathbb R^3$ of two disjoint standard two-spheres joined by an arc, crossed with an interval. Equivalently it is a four-ball with neighborhoods of two disjoint standard proper arcs removed. The barbell diffeomorphism $\beta\in\operatorname{Diff}_{\partial}(B)$ is the point-pushing construction described in the source: loop a middle portion of one deleted arc once around the two-sphere surrounding the other, using the latitude sweep that contracts at its poles and is stationary near the interval endpoints; extend this loop of arcs to an isotopy of the ambient four-ball. Its terminal map fixes the two arc neighborhoods and the outer boundary, and restricts to $\beta$ on $B$. The two sweep orientations give inverse maps. For a smooth embedding $f:B\hookrightarrow S^4$, the implantation is
\[
f_*\beta(x)=
\begin{cases}
f\beta f^{-1}(x),&x\in f(B),\\
x,&x\notin f(B).
\end{cases}
\]
Because $\beta$ is the identity near the boundary, this is a smooth diffeomorphism.
\paragraph{Statement recorded in the supplied source.}
(a) Determine the group $\pi_0(\operatorname{Diff}^+(S^4))$; in particular, is it the trivial group? (b) Does there exist an implantation $f$ with
\[
[f_*\beta]\ne[\mathrm{Id}_{S^4}]\quad\text{in}\quad\pi_0(\operatorname{Diff}^+(S^4))?
\]
The source records two further questions. Let $f_0$ be the winding implantation of its Figure 4 (PDF page 246): the two spheres are separate standard spheres, while the connecting bar travels from the left sphere to link the right sphere once, returns to link the left sphere, then attaches to the right sphere. Use the framing and path represented by that figure. Is $[f_{0*}\beta]$ nontrivial? Finally, is the whole component group generated by implanted barbell classes, and can a finite collection of such classes generate it? Explicitly, is every component a finite product of classes $[f_*\beta]$ and their inverses; can the embeddings $f$ be chosen from one fixed finite collection? The question concerns components, not a homotopy equivalence between the entire diffeomorphism group and $SO(5)$. \sref{2943}{2}
\subsection{Short English statement}
Determine the four-sphere mapping class group, test implanted barbell maps including the specified winding example, and decide whether barbell classes generate the group or a finite collection suffices.
\subsection{Sources}
\begin{description}
\item[\hypertarget{source:2943:1}{S1}] ULAM.AI and the UnsolvedMath Contributors, UnsolvedMath: A Curated Collection of Open Mathematics Problems, record 2943 (KP-4.67). CC BY 4.0. \url{https://huggingface.co/datasets/ulamai/UnsolvedMath}.
\item[\hypertarget{source:2943:2}{S2}] R. İnanç Baykur, Robion C. Kirby and Daniel Ruberman, K3—A New Problem List in Low-Dimensional Topology, Mathematical Surveys and Monographs 295 (2026), author version, Problem 4.67 and Remarks, PDF page 244. \url{https://math.berkeley.edu/sites/default/files/surv-295-ruberman-watermarked-author-pdf.pdf}.
\end{description}
\label{end:2943}
\end{document}
